\documentclass{article}
\usepackage{decomp}
\begin{document}

\decompnote{Full verbatim transcript of US patent 8,109,184 B2 (11 pages: front page, 4 drawing sheets, 6 pages of description and claims in 12 numbered columns), produced tool-first: the text was placed from a Datalab API (accurate mode) OCR of the USPTO image scan made on 2026-10-20 (\texttt{work/datalab/US8109184B2.md}) by the script \texttt{work/build\_main.py}, and proofread by the decompiler against the page images and, word by word, against the text of the same patent on Google Patents (\texttt{source/US8109184B2.google\_patents.txt}); every correction is in the script's correction table and in the README. The source PDF (USPTO full-text image, 11 one-bit CCITT page images of 2560\,\texttimes\,3300 px, i.e. the letter page at 300 ppi, declared as a 2560\,\texttimes\,3300 pt page box) has no text layer; the page images are authoritative. The document is a USPTO patent publication and prints no copyright notice. Left out: the barcode and its number US008109184B2 at the top of the front page, the running heads of the drawing sheets (``U.S. Patent\quad Feb. 7, 2012\quad Sheet N of 4\quad US 8,109,184 B2'') and of the text pages (``US 8,109,184 B2''), the column numbers 1--12 (given in a note at the start of each text page instead) and the printed line numbers 5, 10, \ldots, 65 between the columns. The letter Z of ``Z-axis'' is printed throughout as a Z with a horizontal stroke through it (the patent drafting convention that distinguishes it from the digit 2); it is written as a plain Z here. The reference numerals of the drawings and the figure numbers are printed in bold in the description and the claim numbers in the claims; that bold is kept. The drawing sheets are printed sideways; the figure files are rotated so that the numerals read upright.}

\sourcefile{US8109184B2.pdf}

\pdfpage{1}
(12) \textbf{United States Patent} \textbf{Kong et al.}

(10) \textbf{Patent No.:} \textbf{US 8,109,184 B2}

(45) \textbf{Date of Patent:} \textbf{Feb. 7, 2012}

\section*{(54) METHOD AND DEVICE FOR CUTTING FRESH TISSUE SLICES}

(76) Inventors: \textbf{Jian-Qiang Kong}, Greenville, NC (US); \textbf{George Y Kong}, Greenville, NC (US)

( * ) Notice: Subject to any disclaimer, the term of this patent is extended or adjusted under 35 U.S.C. 154(b) by 395 days.

(21) Appl. No.: \textbf{12/310,519}

(22) PCT Filed: \textbf{May 22, 2007}

(86) PCT No.: \textbf{PCT/US2007/012120}

§ 371 (c)(1), (2), (4) Date: \textbf{Feb. 27, 2009}

(87) PCT Pub. No.: \textbf{WO2008/030279}

PCT Pub. Date: \textbf{Mar. 13, 2008}

\subsection*{(65) Prior Publication Data}

US 2010/0175520 A1 Jul. 15, 2010

\subsection*{Related U.S. Application Data}

(60) Provisional application No. 60/842,386, filed on Sep. 6, 2006.

(51) \textbf{Int. Cl.} \textbf{\textit{G01N 1/06}} (2006.01) \textbf{\textit{B26D 1/45}} (2006.01)

(52) \textbf{U.S. Cl.} \dotfill \textbf{83/13;} 83/646

(58) \textbf{Field of Classification Search} \dotfill 83/13, 915.5, 83/401, 703, 410, 646, 705, 409.2, 563, 613, 83/638

See application file for complete search history.

\subsection*{(56) References Cited}

\subsubsection*{U.S. PATENT DOCUMENTS}

\begin{center}
\begin{tabular}{@{}lllll@{}}
3,828,641 & A & 8/1974 & Sitte & \\
5,181,443 & A * & 1/1993 & Sitte et al. & 83/72\\
5,226,335 & A & 7/1993 & Sitte et al. & \\
5,865,081 & A * & 2/1999 & Myers & 83/149\\
5,906,148 & A & 5/1999 & Aihara et al. & \\
6,209,437 & B1 * & 4/2001 & Izvoztchikov et al. & 83/707\\
6,619,173 & B2 & 9/2003 & Bellasalma & \\
6,651,538 & B2 * & 11/2003 & Tamura et al. & 83/575\\
7,146,895 & B2 & 12/2006 & Kong et al. & \\
\end{tabular}
\end{center}

\subsubsection*{FOREIGN PATENT DOCUMENTS}

\begin{center}
\begin{tabular}{@{}llll@{}}
JP & 57100335 & A & 6/1982\\
WO & 03/093799 & A1 & 11/2003\\
\end{tabular}
\end{center}

\subsubsection*{OTHER PUBLICATIONS}

King, Marcus J., ``Slicing frozen meat with oscillating knife.'' Meat Science, vol. 51, 1999, pp. 251-269.

* cited by examiner

\textit{Primary Examiner} — Ghassem Alie

(74) \textit{Attorney, Agent, or Firm} — Richard C. Litman

\subsection*{(57) ABSTRACT}

The method and device for cutting fresh tissue slices utilizes a vibrating microtome. The vibrating microtome uses a cutting blade that vibrates in a horizontal direction. The microtome has a base, a specimen holder attached to the base, and a diagonal slide carriage and support arm that are supported on the base for movement diagonally with respect to the specimen holder. A vibrating blade support is located on the support arm and moves horizontally with respect to the specimen holder. A cutting blade is mounted on the vibrating blade support and moves in a compound diagonal and vibrating motion with respect to the specimen holder.

\textbf{14 Claims, 4 Drawing Sheets}

\begin{center}
\includegraphics[width=0.8\linewidth,height=0.5\textheight,keepaspectratio]{figures/front_drawing.png}
\end{center}
\alttext{Front page, lower half: the representative drawing of the patent, which is FIG. 1 (sheet 1) reproduced upright and without the ``FIG. 1'' label; the same line drawing with the same 49 reference labels (10, 12, 14, 16, 22, 26, 28, 110, 112, 114, 116, 118, 120 twice, 122, 124, 126, 128, 130, 132, 134, 140, 142, 144, 146, 150, 152, 154, 156, 158 twice, 160, 162, 164, 166, 168, 170, 172, 174 twice, 180, 192 twice, 194, 200, 202, 204, and the arrows B and C). See the description of FIG. 1 on sheet 1 for what every numeral points to. Rendered at 200 ppi from the 300-ppi scan (\texttt{figures/front\_drawing.png}).}

\pdfpage{2}
\begin{center}
\includegraphics[width=0.95\linewidth,height=0.6\textheight,keepaspectratio]{figures/fig1.png}
\end{center}

\noindent\textbf{FIG. 1}
\alttext{Sheet 1 of 4, FIG. 1: a front elevation view of the vibrating sliding-blade microtome, a black-and-white patent line drawing (no shading except the bolt symbols), printed sideways on the sheet and shown here rotated upright; the label ``FIG. 1'' stands at the left middle. Every part carries a reference numeral joined to it by a short wavy leader line; the numerals and the parts they point to, with the description's own names, are as follows. Bottom: a long flat horizontal bar across the whole width is the generally planar base 12 (numeral at the lower left). Rising from the base at the left is the support column 114, drawn as a tall rectangle outline that reaches almost to the top of the drawing; the numeral 114 is at the lower left outside it. Above and around the column is a larger rectangular outline whose top edge carries, at the top centre, a square box: the slide platform drive motor controller 128 (the controller is also the common controller of the two motors). The arrow 10 at the upper right points at the whole microtome; the arrow 14 (upper middle right) points at the microtome assembly, i.e. the sliding and vibrating mechanism; the arrow 16 at the lower right points at the specimen holding assembly (the bath). Inside the column's upper part, an inner rectangle holds the diagonal sliding component 110 (arrow at the lower left of it): the diagonal slide rail 118 is a narrow bar inclined at about 45°, rising from lower right to upper left, carrying two star-shaped fastener bolt symbols 120 at its upper end (numeral 120 at the upper left); a second numeral 120 at the right middle points at a bolt symbol on a short block at the lower right end of the rail. The diagonal slide platform 116 is a wider rectangle lying parallel to and over the rail. Along the lower edge of the platform runs a long narrow inclined bar, the gear tooth rack 122 (numeral below its lower end); the double-headed arrow B, drawn parallel to the rack, shows the diagonal sliding direction. Under the rack a small circle, the drive pinion 124, sits on a large circle, the diagonal slide platform drive motor 126 (both numerals below the circles); a plain line runs from the motor down to the inner rectangle's bottom edge. A small dashed circle at the lower right corner of the platform is the magnet 134. Two small squares are the Hall-Effect sensors: 130 at the left, level with the upper end of the rail, joined by a line to the outer frame, and 132 at the lower right end of the rail, joined by a line downwards. At the right end of the platform a hexagon, the bolt 144, fastens the first, attached end 142 of the support arm 140, a long horizontal rectangle that extends to the right over the bath and ends in its free end 146 (numeral at the right). On the right half of the arm sits the blade support rail 156 (a horizontal bar, numeral at the right below the arm), fastened by the bolts 158 (two star symbols, numerals 158 at the lower left and upper right of the rail), and on it the two blade support carriers 154 (left) and 152 (right), small rectangles that straddle the vibrating blade support 150, a tall vertical plate passing through the arm and hanging below it; the double-headed horizontal arrow C beside its lower part shows the horizontal vibration. The horizontal vibrating component as a whole is pointed at by the arrow 112 (left of the plate). At the upper right, above the arm: the vibrating blade drive motor 160 is the large circle at the far right, the drive wheel 162 the circle drawn inside it, and the eccentrically situated drive pin 164 the small circle on the wheel (numeral 164 above); the connector link 168 is a short bar with rounded ends joining the pin 164 to the blade support drive pin 172 at the top of the plate 150; 166 is the first end of the link (at the pin 164) and 170 its second end (at the pin 172); the two numerals 174 point at the bearing assemblies drawn as concentric circles at both ends of the link. At the bottom of the plate 150 hangs the blade holder, an inverted T: the horizontal dashed bar through the plate's lower end 204 is the blade holder support rod 202, the upper end 200 of the blade holder leg 194 sits on it, the leg runs down to the blade holder bar 192 (a horizontal bar, numeral given twice, left and right), and the thin strip below the bar is the cutting blade 180. Lower right: the specimen receiving bath, a rectangle with a wavy line near its top (the liquid surface), its left side wall 26 (numeral at the left), right side wall 28 and front wall 22 (numerals at the right); inside it a dashed rectangle and, within that, a circle containing an irregular closed outline (the end of the syringe tube with the specimen, not numbered in this view). No scale, no text other than the numerals, letters and the figure label.}

\pdfpage{3}
\begin{center}
\includegraphics[width=0.95\linewidth,height=0.6\textheight,keepaspectratio]{figures/fig2.png}
\end{center}

\noindent\textbf{FIG. 2}
\alttext{Sheet 2 of 4, FIG. 2: a side elevation view, partly in section (cut parts hatched), of the microtome, showing the specimen holding assembly with a specimen being advanced; the label ``FIG. 2'' is at the bottom centre. Orientation after rotating the sheet upright: the bath is at the left, the micrometer drive at the right, the blade at the upper left. Bottom: the base 12, a long flat bar (numeral at the lower left). On it at the left stands the specimen receiving bath (arrow 20 at the upper left): its front wall 22 is the tall hatched wall at the far left, the bath bottom 30 the hatched floor (numeral below it), and the rear wall 24 the hatched wall at the right of the bath, interrupted by the aperture 36 through which the syringe tube passes (numeral 24 at the small hatched block above the tube, 36 below the tube). The saline solution 32 fills the bath: its label is underlined inside the bath and a wavy line from the front wall to the syringe tube marks its surface. To the right of the rear wall stands the syringe support base 40, a large hatched block; its front wall 44 (numeral below the base) faces the bath's rear wall, its rear wall 66 is its right face, and its base aperture 42 is the gap in which the tube lies. The resilient O-ring 38 is drawn as two small circles (the ring in section) above and below the tube, sandwiched between the rear wall 24 and the front wall 44. The syringe tube 34 runs horizontally from inside the bath to the right past the support base: a tube with hatched top and bottom walls; its outlet end 58 is the left end, where the walls curve inwards to the reduced diameter outlet lip 60 (numerals 60 and 58 at the lower left of the tube); its inlet end 62 is the right end (numeral below the tube). The stop 64 is the hatched block on top of the tube just right of the rear wall. Inside the tube at the right is the syringe plunger 50, a large hatched block, with its leading face 52 (left face, inside the tube) and its outer face 92 (right face, outside the tube). Between the leading face and the outlet lip the tube holds the tissue specimen 54 (numeral above, leader to the interior) in the low gel point agarose 56 (numeral below, leader to the wavy-bounded material at the outlet). The arrow 16 (upper middle) points at this specimen holding assembly. Right half: the micrometer drive (arrow 70 at the upper right) in the micrometer drive housing 72, a large rectangle; inside it the micrometer 80 (arrow): the cross-hatched knurled thimble is the micrometer body 84, the graduated sleeve with scale lines leads left into the micrometer plunger 86, which passes out of the housing through the mounting plate 88 (a small plate with two screw symbols on the housing's left face) and ends in its free end 90, a thin pin touching the plunger outer face 92. The support pedestal 82 is the block under the micrometer inside the housing. Under the housing lies the micrometer drive housing carrier plate 74 (a flat plate with a screw hole at each end) on the micrometer drive housing linear slide 76 (a thin rail on the base). Upper left, above the bath: the blade holder (arrow 182) and the cutting blade: the blade holder bar 192 is a hatched bar inclined at about 60° from the vertical, leaning to the upper left, with the horizontal ridge 196 as a rounded bump at its upper end and its lower edge 198 at the bottom; the cutting blade 180 (arrow) is the thin hatched strip glued along the bar's lower face, its blade body 186, its top edge 188 abutting the ridge, its blade side 190 against the bar, and its cutting edge 184 at the tip, which protrudes below the bar's lower edge. A double-headed curved arrow under the tip, labelled D, shows the arc over which the blade's inclination can be varied, and a vertical dashed line from the tip down to the outlet lip marks the cutting plane. No scale.}

\pdfpage{4}
\begin{center}
\includegraphics[width=0.95\linewidth,height=0.6\textheight,keepaspectratio]{figures/fig3.png}
\end{center}

\noindent\textbf{FIG. 3}
\alttext{Sheet 3 of 4, FIG. 3: a top plan view of the micrometer drive assembly, a line drawing; the label ``FIG. 3'' at the left middle. Centre: the micrometer drive housing 72, a rectangle (numeral at the right; the arrow 70 at its upper right corner points at the micrometer drive). Behind the housing, seen above and below it, lies the micrometer drive housing linear slide 76, a long vertical rectangle with four screw heads drawn as circled crosses, two at the top and two at the bottom (numeral 76 at the lower right). Inside the housing at the top is the micrometer drive motor 96, a square box; at its right a small rectangle, the micrometer motor drive gear 98, from which the micrometer drive belt 100 (a long vertical band, drawn with dashed sides where it passes behind the gears) runs down to the micrometer drive gear 102, a larger rectangle at the lower right of the housing. In the lower part of the housing lies the micrometer, drawn horizontally: the cross-hatched knurled thimble, the graduated sleeve with scale lines, and the plunger extending to the left out of the housing through a rectangular plate and ending in a thin pin that touches a hatched block with a curved left face at the far left (the syringe plunger, not numbered in this view). Upper left, outside the housing: the micrometer motor control 104, a square box joined to the motor 96 by a line. At the lower left of the housing a vertical double-headed arrow is labelled ``19''; the description calls this arrow A (the directions in which the housing is slidable along the slide); no part 19 exists. No scale.}

\pdfpage{5}
\begin{center}
\includegraphics[width=0.95\linewidth,height=0.6\textheight,keepaspectratio]{figures/fig4.png}
\end{center}

\noindent\textbf{FIG. 4}
\alttext{Sheet 4 of 4, FIG. 4: a bar graph (the label ``FIG. 4'' below it), two bars on a baseline with a single vertical axis at the left. The y axis is linear with tick marks and labels 0, 2, 4, 6, 8 (the ticks are unevenly spaced in the print, by up to 0.1 unit); its label, printed vertically beside it, is defective: it reads ``Live/'' followed by a small fragment, a gap, and ``l Cell Ratio'' (the word between, Dead, is missing but for its first and last strokes; the description calls the quantity the live/dead cell ratio). The x axis has no labels or ticks. Left bar, open (white), reference numeral 206 with a wavy leader to its right edge: the slice viability (live/dead cell ratio) of slices prepared with the vibrating, sliding-blade microtome of the invention, about 6.2 high, with an error bar rising from the bar top to a cap at about 6.9 and an asterisk above the cap (not explained in the patent). Right bar, hatched with diagonal lines, numeral 208 with a leader to its right edge: the ratio obtained with a prior vibrating microtome, about 1.2 high, error bar to about 1.4. No sample sizes, statistics or error-bar definition are given. Values digitized in \texttt{figures/fig4.csv} (uncertainty about ±0.1).}
\figuredata{FIG. 4}{figures/fig4.csv}{figures/fig4.py}{digitized (tier 3) from the 300-ppi bilevel scan of sheet 4 with figures/fig4\_extract.py, axis calibrated on the five printed ticks; the ticks are unevenly printed, so about ±0.1 ratio units}

\pdfpage{6}
\decompnote{PDF page 6 prints columns 1 (left) and 2 (right); column 2 begins at ``reduce this Z-axis movement''.}
\section*{METHOD AND DEVICE FOR CUTTING FRESH TISSUE SLICES}

\subsection*{CROSS-REFERENCE TO RELATED APPLICATIONS}

This application claims priority, under 35 USC 120, to U.S. provisional patent application No. 60/842,386, filed Sep. 6, 2006, the disclosure of which is expressly incorporated herein by reference.

\subsection*{FIELD OF THE INVENTION}

The present invention is directed generally to a vibrating microtome. More specifically, the present invention is directed to a vibrating microtome that utilizes a blade which is movable in two directions. Most specifically, the present invention is directed to a vibrating microtome, and to a method of its use in the preparation of live or fixed tissue slices. The tissue specimens are embedded inside a specimen syringe. A slice thickness is selected using a stepper motor to drive a micrometer. A cutting blade is caused to slide diagonally, and to vibrate horizontally to slice very thin layers or slices off the specimen which is caused to extend from the specimen syringe. The vibrating and sliding blade microtome, in conjunction with the tissue embedding and slice thickness determination provides a much improved tissue slice quality.

\subsection*{BACKGROUND OF THE INVENTION}

Microtomes are generally well-known and are typically used in biomedical experiments to prepare thin slices of tissue that can then be subjected to microscopic studies. The tissue to be sliced is often embedded in paraffin or is hardened by freezing so that it can be sliced as thinly as possible, typically through the use of a microtome. Live tissues and freshly fixed tissues are typically too soft for this type of cutting and it is thus difficult to obtain thin slices of such live tissues and freshly fixed tissues. One attempted solution has been the provision of a vibrating blade microtome in which the cutting blade has been vibrated horizontally with respect to the typically vertical or perpendicular movement of the cutting blade edge. Such vibrating microtomes are specifically intended to cut soft tissue slices without the need for tissue embedding or freezing. The horizontal oscillation of the blade, or its oscillation is a direction which is generally parallel to the longitudinal axis of the cutting edge, is provided as a high frequency vibration that has been more effective in slicing soft tissue samples than were prior cut devices. The resultant slices, cut by the vibrating microtomes, have had better slice quality than have slices cut by direct blade chopping.

In all current vibrating microtomes, the blade vibrates generally parallel with respect to the longitudinal axis of the cutting blade. Typically this means that the blade vibrates in a generally horizontal, side to side direction, at a high frequency. An amplitude of such a vibrating stroke must be at least 0.6 mm. If the stroke amplitude is less than this amount, the blade will not cut through the live tissue. Instead, the live tissue or the freshly fixed specimen will merely be compressed and will push back against the horizontally vibrating blade.

Vibrating microtomes have been noted recently as having undesirable vibrations of the cutting blade in the Z-axis direction; i.e. the direction perpendicular to both the X-axis blade vibration direction and the Y-axis cutting direction. Efforts have been made in prior vibrating microtomes to reduce this Z-axis movement of the cutting blade using a so-called ``zero —'' technique. While the blade movement is most often tested and quantified in test conditions, during vibrating of the blade but without cutting a tissue slice, the Z-axis movement or vibration of the cutting blade in operating conditions has been largely ignored.

Double edged razor blades are widely used in vibrating microtomes for cutting tissue slices. Such a razor blade is broken or separated into two single edge blade parts by being severed or snapped along its longitudinal axis prior to use. The now half blade, with its single cutting edge is clamped, generally at its middle, along its longitudinal axis in the blade holder of a currently available vibrating microtome and is used to cut tissue slices. The thin razor blade is apt to be bent or deflected relative to the blade holder, particularly in the portion of the blade that is unsupported by the blade holder of the vibrating microtome. The razor blade is typically mounted on the blade holder and is inclined at an angle of 15-35 degrees, with respect to the surface of the tissue specimen. During slicing of the specimen, the specimen will itself tend to pull and to bend the blade toward the specimen. Such a bending of the blade, during slicing of the specimens, causes the formation of chatter marks in the cut or sliced specimen because of blade bending, or chatter, in the Z-axis direction. Resistance of the tissue specimen to being sliced, and the inherent thin blade structure, together with its being mounted in a somewhat flexible fashion, all contribute to flexure of the blade in the Z-axis and to the resultant formation of tissue damage by the infliction of chatter marks on the surface of the cut tissue specimens.

In current vibrating machines, the edge of the cutting blade vibrates on a horizontal cutting plane. Such a vibrating motion requires a long vibrating arm so that the cutting blade will be located in the middle of the tissue bath container. This relatively long vibrating arm is heavy and is difficult to support so that it will be able to vibrate at the desired speed but will not wobble.

In U.S. Pat. No. 7,146,895, to Kong et al., the applicants in the subject application, there is disclosed a sliding blade microtome. As described in that patent, the disclosure of which is expressly incorporated herein by reference, a diagonally sliding cutting blade is usable to slice thin specimens of live or dead tissue, which specimens have been embedded in a gel and have been extended from the open end of a specimen-holding syringe, by the use of a sliding blade microtome machine. The sliding blade machine disclosed in this prior Kong patent can cut soft tissue such as brain cerebrum into slices. However, this prior device still has several limitations. A first of these is that in connection with high fiber containing tissue, which thus tends to be difficult to cut, such as brain stem or cerebellum, direct sliding cannot cut through the tough tissue thus resulting in low quality tissue slices. Also, the thinnest slices that can be cut with the sliding blade machine described in the above-identified prior Kong patent have a thickness of about 70 to 80 μm. This thickness of cut precludes the sliding blade machine from being used in histology studies which typically require slice thicknesses between 10 and 40 μm.

It will thus be seen that a need exists for a device for cutting fresh tissue slices and for a method of using such a device that overcomes the limitations of the presently existing procedures and apparatus. The method and device for cutting fresh tissue slices in accordance with the present invention provides such a device and method and is a substantial advance in the art.

\subsection*{SUMMARY OF THE INVENTION}

It is an object of the present invention to provide a device for use in cutting fresh tissue slices.

\pdfpage{7}
\decompnote{PDF page 7 prints columns 3 (left) and 4 (right); column 4 begins at ``vibrating unit can move freely''.}
Another object of the present invention is to provide a vibrating blade microtome.

A further object of the present invention is to provide a vibrating sliding-blade microtome.

Yet another object of the present invention is to provide a method for using a vibrating, sliding-blade microtome to reduce tissue damage.

Still a further object of the present invention is to provide a vibrating, sliding-blade microtome that eliminates chattering in a Z-axis direction.

Even yet another object of the present invention is to provide a vibrating, sliding-blade microtome which reduces mechanical trauma on tissue slices.

The method and device for cutting fresh tissue slices, in accordance with the present invention utilizes a vibrating microtome, in which the cutting blade vibrates in a generally horizontal direction, with a sliding blade mechanism that causes the blade to also travel in a diagonal X-Y axis direction. The vibrating sliding-blade microtome in accordance with the present invention also adheres the cutting blade to one surface of a blade holder in a procedure that substantially reduces bending and resultant chattering of the blade in the Z-axis direction. The method and device in accordance with the present invention is particularly intended for use in the slicing of live or fixed tissue slices which are not hard or rigid. Such specimens are not embedded in paraffin and are not frozen. The vibrating sliding-blade microtome in accordance with the present invention cuts specimen slices that have substantially higher numbers of live/dead cell ratios than are capable of being obtained with prior art devices. Such a higher number of live/dead cell ratio is an indication of slice viability.

The vibrating, sliding-blade microtome in accordance with the present invention combines a transverse or back and forth, generally horizontal vibrating motion with a diagonal downward sliding motion. The resultant compound cutting blade movement has a substantial effect on the cutting efficiency of the vibrating sliding-blade microtome of the present invention. Whereas, in the prior art devices, the blade vibration amplitude had to be at least 0.6 mm, in the subject device, that amplitude can be as little as 0.1 mm. The specimen being cut is not compressed by the blade vibrating movement so that the cutting efficiency of the blades is increased, and the mechanical stress imparted to the specimen is reduced.

In contrast with prior art devices, the razor blade, or more typically, the half of a broken double-edged razor blade is glued onto a one piece blade holder. This allows the blade to be securely attached to the blade holder all along its length and at one side of the blade holder. The glue can extend to the edge of the blade holder and thus can extend to within 0.5 mm of the cutting edge of the razor blade. This forms a very rigid, uniform support for the thin razor blade and essentially eliminates blade edge wobble. In clamped blades, there are typically one or more clamping screws or the like. The blade is unevenly stressed and is subjected to blade edge wobble. The elimination of blade edge wobble by the use of the blade mounting through the use of adhesives of the present invention, in combination with the vibrating, sliding blade movement of the blade, essentially eliminates Z-axis chatter, which is prevalent in prior devices.

In the present invention, the cutting plane is generally vertical, in contrast to the prior art devices, in which the cutting plane is typically horizontal. This means that there are not required the long vibrating arms which were common in prior devices. The vibrating arm can essentially be as short as the cutting stroke and is thus very light weight. The vibrating unit is mounted on a preloaded linear bearing so that the vibrating unit can move freely in the vibrating, horizontal or X direction with no component of movement, by itself in either of the Y or Z axis directions. Very low wattage is required to move the vibrating element as compared to the higher levels of power required by prior vibrating microtomes.

The specimen to be sliced by the vibrating, sliding-blade microtome of the present invention is held in a specimen holder, as disclosed in detail in the prior Kong et al. U.S. Pat. No. 7,146,895. The combination of the specimen support structure and process, as described in the prior Kong et al. patent, together with the utilization of the vibrating, sliding-blade microtome of the present invention results in the preparation of extremely thin, sharply and accurately cut tissue specimens, particularly live tissue specimens. The slice viability, which is an important factor in the evaluation of a fresh, live tissue slice, is much greater than has been accomplishable using prior art devices. The live/dead cell ratio, which is a measurement of slice viability, is much greater in the present device, than has been the case when using traditional vibrating microtomes.

The method and device for cutting fresh tissue slices, in accordance with the present invention, overcomes a number of the significant limitations of the prior art devices. It is a substantial advance in the art.

\subsection*{BRIEF DESCRIPTION OF THE DRAWINGS}

A full and complete understanding of the method and device for cutting fresh tissue slices, in accordance with the present invention may be had by referring to the detailed description of the preferred embodiment, as is set forth subsequently, and as illustrated in the accompanying drawings, in which:

FIG. \textbf{1} is a front elevation view of a vibrating sliding-blade microtome in accordance with the present invention;

FIG. \textbf{2} is a side elevation view, partly in section, of the vibrating, sliding-blade microtome of FIG. \textbf{1} and showing a specimen being advanced in the specimen holder prior to its being sliced;

FIG. \textbf{3} is a top plan view of a portion of the vibrating, sliding-blade microtome of the present invention and showing a micrometer drive assembly; and

FIG. \textbf{4} is a bar graph depiction of live/dead cell ratios obtainable by use of the present invention and compared with the results obtainable using a prior art device.

\subsection*{DESCRIPTION OF THE PREFERRED EMBODIMENT}

Referring initially to FIG. \textbf{1}, there may be seen, generally at \textbf{10} a preferred embodiment of a vibrating sliding-blade microtome in accordance with the present invention. A generally planar base \textbf{12} supports a microtome assembly, generally at \textbf{14}, and a specimen holding assembly, generally at \textbf{16}. As will be discussed in detail shortly, the specimen holding assembly, generally at \textbf{16}, is utilized to support and hold a tissue specimen so that a thin slice of the tissue can be cut from the specimen for further analysis. The specimen holding assembly, generally at \textbf{16}, is similar, in overall structure and operation, to the specimen holding assembly described and depicted in the prior U.S. Pat. No. 7,146,895 in the names of the present applicants. As indicated above, the entire disclosure of that patent is expressly incorporated herein be reference.

\decompnote{PDF page 8 begins inside the next paragraph (at the inline page marker); it prints columns 5 (left) and 6 (right); column 6 begins at ``electric stepper motor.''.}
As may be seen by referring to FIG. \textbf{2}, taken in conjunction with FIG. \textbf{1}, the specimen holding assembly, generally at \textbf{16}, \pdfpage{8}includes a specimen receiving bath, generally at \textbf{20} which includes a bath front wall \textbf{22}, a bath rear wall \textbf{24}, left and right bath side walls \textbf{26} and \textbf{28}, as seen in FIG. \textbf{1} and a bath bottom \textbf{30}. The specimen receiving bath \textbf{20} thus is a fluid-tight receptacle into which can be placed a suitable saline solution \textbf{32} or any other fluid whose purpose will be to maintain the viability of a specimen slice that is deposited in the bath \textbf{20} by the action of the microtome assembly, generally at \textbf{14}, as will be discussed subsequently.

Referring again to FIG. \textbf{2}, a syringe tube \textbf{34} is insertable into an aperture \textbf{36} in the rear wall \textbf{24} of the specimen receiving bath, generally at \textbf{20}. A resilient O-ring \textbf{38} surrounds the bath aperture \textbf{36} and will prevent leakage of the saline solution \textbf{32} or other liquid out of the specimen receiving bath \textbf{20} while the syringe tube \textbf{34} is inserted through the aperture \textbf{36}. It will be apparent that the saline solution or other liquid solution \textbf{32} will be added to the bath only after the syringe tube \textbf{34} is in place and will be drained before the syringe tube \textbf{34} is removed. A syringe support base \textbf{40} is situated spaced behind, and generally parallel to the rear wall \textbf{24} of the specimen receiving bath. That base is provided with a base aperture \textbf{42} which is aligned with the bath rear wall aperture \textbf{36}. As may also be seen in FIG. \textbf{2}, the resilient O-ring \textbf{38} is sandwiched between the rear wall \textbf{24} of the specimen receiving bath \textbf{20} and a front wall \textbf{44} of the syringe support base \textbf{40}.

The syringe tube \textbf{34} is a generally cylindrical tube that is sized to receive a cooperating syringe plunger \textbf{50}. A leading face \textbf{52} of the syringe plunger \textbf{50} is adapted to have attached thereto a tissue specimen \textbf{54} which can be attached to the syringe leading face \textbf{52} by a suitable instantaneous adhesive or the like. Tissue specimen \textbf{54} will typically be encased in, or embedded in a low gel point agarose \textbf{56}. As may be seen in FIG. \textbf{2}, the syringe tube \textbf{34} has an outlet end \textbf{58} with a reduced diameter outlet lip \textbf{60} whose purpose is to provide some overcomable resistance to the passage of the tissue specimen \textbf{54} and its agarose encasement \textbf{56} out of the outlet end \textbf{58} of the syringe tube \textbf{34}. Clearly, the reduced diameter outlet lip \textbf{60} does not provide a complete barrier to the passage of the tissue specimen \textbf{54} and its embedding agarose \textbf{56} out of the syringe tube outlet end \textbf{58}.

The syringe tube \textbf{34} is also structured with an inlet end \textbf{62} which, in the present invention is sized to slidingly receive the syringe plunger \textbf{50}. A stop \textbf{64} is attached to the syringe tube \textbf{34} intermediate the syringe tube ends \textbf{58} and \textbf{62} respectively. Stop \textbf{64} is positioned so that it will abut a rear wall \textbf{66} of the syringe support base \textbf{40}. Stop \textbf{64} thus functions to ensure that the syringe tube \textbf{34} will extend into the saline solution \textbf{32} or other liquid in the specimen receiving bath \textbf{20} to a pre-determined, repeatable depth or position. This will aid in the aligning of the vibrating sliding-blade microtome with the outlet end \textbf{58} of the syringe tube \textbf{34}.

A micrometer drive, generally at \textbf{70}, is, as may be seen in FIGS. \textbf{2} and \textbf{3}, situated in a micrometer drive housing, generally at \textbf{72}. The micrometer drive housing \textbf{72} is, in turn, secured atop a micrometer drive housing carrier plate \textbf{74} which is slidable, in the directions indicated by the arrow A in FIG. \textbf{3}, along a micrometer drive housing linear slide \textbf{76} which is, in turn, secured to the base \textbf{12} by any suitable expedient. A micrometer \textbf{80} is situated atop a support pedestal \textbf{82} in the micrometer drive housing \textbf{72} and has a micrometer body \textbf{84} which carries a micrometer plunger \textbf{86} which passes through the micrometer drive housing \textbf{72} and through a mounting plate \textbf{88} attached to the micrometer drive housing \textbf{72}. A free end \textbf{90} of the micrometer plunger \textbf{86} is engagable with a syringe plunger outer face \textbf{92}.

A micrometer drive motor \textbf{96} is situated in the micrometer drive housing \textbf{72} and is preferably provided as a suitable electric stepper motor. Micrometer drive motor \textbf{96} has a micrometer motor drive gear \textbf{98} which supports a micrometer drive belt \textbf{100} that extends from the micrometer motor drive gear \textbf{98} to a micrometer drive gear \textbf{102}. Both of these drive gears \textbf{98} and \textbf{102} are externally toothed and engage interior teeth or cogs on the drive belt \textbf{100}. The micrometer drive motor \textbf{96} is very precisely controlled by a micrometer motor control \textbf{104}. That control can be actuated to rotate the micrometer drive stepper motor \textbf{96} a finite amount so that the micrometer plunger \textbf{86}, and specifically its free end \textbf{90} will engage the syringe plunger \textbf{50} and can advance it in the syringe tube \textbf{34} a defined amount. Such a defined amount of forward travel of the syringe plunger \textbf{50} will extrude a set length of the specimen \textbf{54} and embedding agarose \textbf{56} out of the reduced diameter outlet lip \textbf{60} of the syringe tube \textbf{34} so that it can be sliced off by the vibrating sliding-blade microtome and deposited in the saline solution \textbf{32} or the other suitable solution in the specimen receiving bath \textbf{20} where it can be removed for further processing. When it is necessary to remove the syringe plunger \textbf{50} out of the rear, inlet end \textbf{62} of the syringe tube \textbf{34}, the micrometer plunger \textbf{86} and its free end \textbf{90} can be retracted toward the micrometer drive housing \textbf{72} by reversal of the micrometer drive motor \textbf{96}. Once the micrometer plunger \textbf{86} has been so retracted, the micrometer drive housing \textbf{72} can be slid laterally on the micrometer drive housing slide \textbf{76}. If desired, suitable clamps, which are not specifically shown, can be used to hold the micrometer drive housing \textbf{72} in its operational position, and can be loosened to allow it to slide laterally to its non-use position.

Referring now again to FIG. \textbf{1}, the vibrating sliding-blade microtome assembly, generally at \textbf{14}, in accordance with the present invention, will be described in detail. It will be quite evident to those of skill in the art that the purpose of the microtome is to slice thin uniform viable tissue samples from the tissue specimen \textbf{54} embedded in the low gel point agarose \textbf{56} and extruded or forced out of the syringe outlet end \textbf{58} and restricted by the reduced diameter outlet lip \textbf{60}. The incremental advance of that tissue specimen \textbf{54} is accomplished by operation of the micrometer stepper drive motor \textbf{96}, under the control of the micrometer motor controller \textbf{104}. The proper embedding of the tissue specimen \textbf{54} in the low gel point agarose \textbf{56} and the cooperation of the syringe tube \textbf{34}, the reduced diameter outlet tip \textbf{60} and the provision and control of the micrometer stepper drive motor \textbf{96} will set up the tissue for its being sliced. However, if the microtome blade is not able to perform an effective cut, the accurate tissue specimen set-up and positioning is of no consequence. In the prior devices, as discussed above, the microtome blade has tended to chatter and to bend or deflect, has tended to compress the tissue instead of cutting it and has generally failed to provide very thin, reproducible tissue specimen slices of live or soft tissue. The vibrating sliding-blade microtome of the subject invention overcomes those limitations.

As may be seen most clearly in FIG. \textbf{1}, the vibrating sliding-blade microtome, generally at \textbf{14}, is comprised essentially of two distinct yet interacting components. The first of these is a diagonal sliding component, generally at \textbf{110}. The second is a horizontal vibrating component, generally at \textbf{112}. While the two components have separate features, they combine to much more effectively slice the tissue specimens \textbf{54}.

\decompnote{PDF page 9 begins inside the next paragraph (at the inline page marker); it prints columns 7 (left) and 8 (right); column 8 begins at ``160 is situated atop the support arm''.}
A support column \textbf{114} extends upwardly from the base \textbf{12} of the vibrating sliding-blade microtome, generally at \textbf{10}. This column \textbf{114} is depicted in FIG. \textbf{1} as being generally rectangular but need not have any specific shape so long as it has sufficient structural rigidity to support the two cooperating components \textbf{110} and \textbf{112}. A diagonal slide platform \textbf{116} is supported for diagonal movement along a diagonal slide rail \pdfpage{9}\textbf{118} which is, in turn, secured to the support column \textbf{114} by suitable fasteners, such as the fastener bolts \textbf{120} depicted schematically in FIG. \textbf{1}. It will be understood that the diagonal slide platform \textbf{116} and the diagonal slide rail \textbf{118} will preferably have friction reducing components, such as linear bearings, low friction slide surfaces or the like so that the platform \textbf{116} will be able to slide easily along the slide rail \textbf{118}. It will also be understood that the angle of inclination of the diagonal slide rail \textbf{118}, with respect to a horizontal line, could be varied such as, for example, by providing arcuate slots in the support column \textbf{114} and having the fastener bolts \textbf{120} be received in those arcuate slots. This angle of inclination is preferably between 20° and 80°. Other angle adjusting expedients are within the purview of the subject invention.

A gear tooth rack \textbf{122} is attached to the lower edge of the diagonal slide platform \textbf{116} and is in gear mesh with a cooperating drive pinion \textbf{124}. That drive pinion is driven by a diagonal slide platform drive motor \textbf{126} which is preferably a stepper motor whose revolutions are accurally controllable by a slide platform drive motor controller \textbf{128}. A pair of Hall-Effect sensors \textbf{130}, \textbf{132} are situated adjacent first and second ends of the diagonal slide rail \textbf{118}. A suitable magnet \textbf{134} is affixed to the diagonal slide platform \textbf{116}. In operation, readings from the two Hall Effect sensors \textbf{130}, \textbf{132} are sent to the diagonal slide platform drive motor controller \textbf{128} and are usable to very accurately position the diagonal slide platform \textbf{116}. It will be understood that other similar control schemes could be used to accurately determine and control the position of the diagonal slide platform \textbf{116} along the diagonal slide rail \textbf{118}. It is important the platform \textbf{116} be capable of sliding in a diagonal direction, with respect to X and Y axes, as indicated by arrow B in FIG. \textbf{1}, and that its positioning and magnitude of slide travel be capable of being accurately controlled.

A support arm \textbf{140} is secured, at a first, attached end \textbf{142}, to the diagonal slide platform \textbf{116} by a suitable fastener, such as a bolt \textbf{144}. It will be understood that the bolt \textbf{144} is a schematic representation of any suitable fastening arrangement that could be used to attach the support arm \textbf{140} to the diagonal slide platform \textbf{116}. Since it is preferable that the support arm \textbf{140} be maintained generally horizontally, the fastening arrangement \textbf{144} could be adjustable so that the support arm \textbf{140} would remain horizontal with varying inclinations of the diagonal slide rail \textbf{118}. A second, free end \textbf{146} of the support arm \textbf{140} extends, in a cantilever fashion, over the specimen holding assembly, generally at \textbf{16}. It will be understood that the support arm \textbf{140} is depicted in the somewhat schematic illustration of the subject invention in FIG. \textbf{1}, as a generally rectangular member. In actuality, the support arm \textbf{140} could take any number of shapes so long as it has sufficient structural rigidity to accomplish its intended task, as will now be discussed.

A vibrating blade support \textbf{150} is carried by a pair of spaced blade support carriers \textbf{152}, \textbf{154} which are slidable horizontally along a blade support rail \textbf{156}. That blade support rail \textbf{156} is, in turn, bolted or otherwise fastened to the support arm \textbf{140} intermediate its ends \textbf{142} and \textbf{146}, as indicated schematically by bolts \textbf{158}. As was discussed in connection with the diagonal slide platform \textbf{116} and the diagonal slide rail \textbf{118}, the blade support carriers \textbf{152} and \textbf{154}, which may be a unified blade support carriage, and the blade support rail \textbf{156} may be provided with friction reducing elements, such as linear bearings, allow friction slide or other friction reducing components and the like. The result is that the vibrating blade support \textbf{150} is supported so that it will move or reciprocate very freely and easily in the horizontal or X-axis direction, as indicated by arrow C in FIG. \textbf{1}. A vibrating blade drive motor \textbf{160} is situated atop the support arm \textbf{140} generally adjacent the free, unsupported second end \textbf{146} of that support arm \textbf{140}. That vibrating blade drive motor \textbf{160} is typically also an electric motor and is controlled by the same motor controller \textbf{128} which also controls the diagonal slide platform drive motor \textbf{126}. Coordination between the two motors \textbf{126} and \textbf{160} is thus assured through the provision of the common controller \textbf{128}. A drive wheel \textbf{162} is driven by this motor \textbf{160} and is provided with an eccentrically situated drive pin \textbf{164}. That drive pin is received in a first end \textbf{166} of a resilient or spring steel connector link \textbf{168}. A second end \textbf{170} of this connector link \textbf{168} is carried by a blade support drive pin \textbf{172}. As indicated schematically in FIG. \textbf{1} both ends \textbf{166} and \textbf{170} of the resilient or spring steel connector link \textbf{168} may be provided with suitable bearing assemblies, generally at \textbf{174}, which bearing assemblies aid in the support and connection of the connector link \textbf{168}.

Actuation of the vibrating blade drive motor \textbf{160} will cause the drive wheel \textbf{162} to rotate at a relatively high rate of speed. The eccentric location of the drive pin \textbf{164} will cause the connector link \textbf{168} to reciprocate. This reciprocatory motion of the connector link \textbf{168} is translated into a horizontal vibration of the vibrating blade support \textbf{150} in the direction indicated by arrow C in FIG. \textbf{1}.

A cutting blade, generally at \textbf{180}, is attached to a blade holder, generally at \textbf{182}, as seen in FIG. \textbf{1}, and as is also seen in greater detail in FIG. \textbf{2}. The cutting blade \textbf{180} has an elongated blade edge \textbf{184} and a blade body \textbf{186}. Typically, the cutting blade \textbf{180} is one half of a double-edged razor blade. The blade body \textbf{186} has a top edge \textbf{188} and a blade side \textbf{190}. The blade holder \textbf{182} is generally in the shape of an inverted T, as seen in FIG. \textbf{1} and includes a blade holder bar \textbf{192} which is attached, generally at its mid-point to a blade holder leg \textbf{194}. As seen in FIG. \textbf{2}, the blade holder bar \textbf{192} has a horizontal ridge \textbf{196} which acts as an abutment which will be engaged by the upper edge \textbf{188} of the blade body \textbf{186}. The blade cutting edge \textbf{184} extends below a lower edge \textbf{198} of the blade holder bar \textbf{192}. In use, the cutting blade \textbf{180} is attachable to the blade holder bar \textbf{192} by the use of a suitable adhesive which will be applied to the entire length of the blade holder bar \textbf{192}, below the ridge \textbf{196} thereof and to the cutting blade side \textbf{190}. The cutting blade \textbf{180} will be adhesively attached to the blade holder bar \textbf{192} so that its cutting edge \textbf{184} protrudes below the blade holder bar lower edge \textbf{198} by a distance of approximately 1 mm. This support of the cutting blade \textbf{180} by the cutting blade holder \textbf{182} prevents the cutting blade \textbf{180} from bending or flexing, as typically happens with blades attached to holders only in several spaced locations.

Referring again to FIG. \textbf{1}, an upper end \textbf{200} of the blade holder leg \textbf{194} is pivotally supported by a blade holder support rod \textbf{202}. That blade holder support rod \textbf{202} is situated adjacent a lower end \textbf{204} of the vibrating blade support \textbf{150}. A suitable securement assembly, such as a clamping screw, a set screw, or the like is provided, either at the juncture of the blade holder leg \textbf{194} and the rod \textbf{202}, or at the juncture of the rod \textbf{202} and the lower end \textbf{204} of the vibrating blade support \textbf{150} so that an angle of inclination of the cutting blade \textbf{180} can be varied over an arc D, with respect to the end of the specimen to be sliced, as depicted in FIG. \textbf{2}.

\decompnote{PDF page 10 begins inside the next paragraph (at the inline page marker); it prints columns 9 (left) and 10 (right); column 10 begins at ``a diagonal slide carriage on said base (the third element of claim 1)''.}
In operation of the vibrating, sliding-blade microtome, in accordance with the present invention, a tissue specimen \textbf{54} is secured to the syringe leading face \textbf{52} and is embedded in a suitable low gel point agarose \textbf{56}. The syringe plunger \textbf{50} is inserted into the syringe barrel or tube \textbf{34}, from the inlet end \textbf{62}. The now assembled syringe is inserted into the aligned apertures of the bath rear wall \textbf{24} and of the syringe support base \textbf{40}. Once this has been done, the specimen receiving bath \pdfpage{10}\textbf{20} can be filled with a suitable saline solution \textbf{32} or other similar sample receiving medium. The micrometer drive \textbf{70}, which had been slid out of the way or out of alignment with the aligned apertures \textbf{36} and \textbf{42}, is now slid along the micrometer drive housing linear slide \textbf{76} until the micrometer plunger \textbf{86}, and specifically its free end \textbf{90}, is aligned with the syringe plunger outer face \textbf{92}. At this point, the micrometer drive motor \textbf{96} with its controller \textbf{104}, which is in connection with the controller \textbf{128} for the diagonal sliding platform drive motor \textbf{126} and the vibrating blade drive motor \textbf{160}, is actuated to advance the tissue specimen \textbf{54} out through the reduced diameter outlet lip \textbf{60} of the syringe tube \textbf{34}.

The diagonal slide platform \textbf{116} is caused to move down and to the right, as seen in FIG. \textbf{1}, in the direction of the arrow B. At the same time, the cutting blade \textbf{180} is caused to vibrate horizontally in the direction of the arrow C. The compound sliding and vibrating motion of the cutting blade \textbf{180} is effective to slice a first tissue slice off the tissue specimen \textbf{54}. That first slice will typically be discarded because its outer surface, the surface not newly formed by the blade \textbf{180}, will not be even. Subsequent slices, which are obtained by the steps of advancing the syringe plunger \textbf{50} and by the vibrating and diagonally sliding movement of the cutting blade \textbf{180}, are much more uniform in thickness and can be retrieved from the saline solution \textbf{32} for further processing.

In the vibrating and sliding-blade microtome in accordance with the present invention, the cycle time is substantially reduced from that which is possible using only a vibrating blade microtome. In addition, the slice thickness is much more uniform than has been possible in prior devices. The Hall Effect sensors \textbf{130} and \textbf{132} limit the travel of the diagonal slide platform \textbf{116} so that there is no unnecessary platform travel. The short vibrating strokes imparted to the vibrating blade support \textbf{150} are very efficient and again reduce specimen cutting or slicing time. The cutting edge \textbf{184} of the cutting blade \textbf{180} moves in a plane, indicated by dashed lines in FIG. \textbf{2}, which is parallel to the plane of the decreased diameter lip \textbf{60} of the syringe tube \textbf{34}, the specimen slices are thus sliced clearly, quickly and with an accurately and controllable thickness.

One effective measure of the viability of a specimen slice is a ratio of live to dead cells. The higher the ratio; i.e. the greater number of live cells there are for each dead cell, is an excellent measure of slice viability. As may be seen in FIG. \textbf{4}, the slice viability \textbf{206} of slices prepared using the vibrating, sliding-blade microtome in accordance with the present invention, as a ratio of live/dead cells, is much better than a ratio \textbf{208} obtained using a prior vibrating microtome. The vibrating, sliding-blade microtome, in accordance with the present invention is thus clearly far superior to the products which are currently available.

While a preferred embodiment of a method and device for cutting fresh tissue slices, in accordance with the present invention has been set forth fully and completely hereinabove, it will be apparent to one of skill in the art that various changes in, for example, the nature of the specimen slice receiving bath, the specifics of the several drive motors and their controllers, the specific adhesive used to secure the cutting blade to the blade holder, and the like could be made without departing from the true spirit and scope of the present invention, which is accordingly to be limited only by the following claims.

\noindent What is claimed is:
\begin{hangpars}

\textbf{1}. A method for cutting a slice of a specimen tissue comprising the steps of:

\hspace*{1.5em}providing a base and a specimen holder attachable to said base, said specimen holder including:

\hspace*{1.5em}a diagonal slide carriage on said base, said diagonal slide carriage being supported for movement diagonally with respect to said specimen holder, wherein the diagonal slide carriage includes a movable diagonal slide platform and a diagonal slide rail and further includes a diagonal slide platform drive motor;

\hspace*{1.5em}a support arm on said diagonal slide carriage and movable with said diagonal slide carriage;

\hspace*{1.5em}a vibrating blade support on said support arm and movable horizontally with respect to said specimen holder; and

\hspace*{1.5em}a specimen cutting blade on said vibrating blade support, said specimen cutting blade being movable in a compound diagonal and vibrating movement with respect to said specimen holder;

\hspace*{1.5em}attaching a specimen to be sliced to said specimen holder;

\hspace*{1.5em}supporting said specimen cutting blade for horizontal vibrating movement and for diagonal sliding movement with respect to said specimen holder;

\hspace*{1.5em}extending a free end of said specimen from said specimen holder; and

\hspace*{1.5em}severing a slice of said specimen by vibrating and diagonally sliding said specimen cutting blade with respect to said specimen.

\textbf{2}. The method of claim \textbf{1} further including providing a vibrating blade drive motor and using said vibrating blade drive motor for effecting said horizontal vibrating movement of said specimen cutting blade.

\textbf{3}. A device for cutting specimen slices comprising:

\hspace*{1.5em}a base;

\hspace*{1.5em}a specimen holder attachable to said base;

\hspace*{1.5em}a diagonal slide carriage on said base, said diagonal slide carriage being supported for movement diagonally with respect to said specimen holder, wherein the diagonal slide carriage includes a movable diagonal slide platform and a diagonal slide rail and further includes a diagonal slide platform drive motor;

\hspace*{1.5em}a support arm on said diagonal slide carriage and movable with said diagonal slide carriage;

\hspace*{1.5em}a vibrating blade support on said support arm and movable horizontally with respect to said specimen holder; and

\hspace*{1.5em}a cutting blade on said vibrating blade support, said cutting blade being movable in a compound diagonal and vibrating movement with respect to said specimen holder.

\textbf{4}. The device of claim \textbf{3} further including a rack on said diagonal slide platform and a coacting pinion on said diagonal slide platform drive motor.

\textbf{5}. The device of claim \textbf{3} further including a motor controller for said diagonal slide platform drive motor.

\textbf{6}. The device of claim \textbf{5} further including spaced Hall Effect sensors on said base and a magnet on said diagonal slide platform, said Hall Effect sensors being connected to said motor controller.

\textbf{7}. The device of claim \textbf{3} further including a vibrating support drive motor on said support arm.

\textbf{8}. The device of claim \textbf{7} further including an eccentric drive on said vibrating blade support drive motor.

\textbf{9}. The device of claim \textbf{7} further including a connecting link between said vibrating blade support drive motor eccentric drive and said vibrating blade support.

\textbf{10}. The device of claim \textbf{9} wherein said connecting link is spring steel.

\textbf{11}. The device of claim \textbf{3} further including a vibrating blade support carriage connecting said vibrating blade support and said support arm.

\pdfpage{11}
\decompnote{PDF page 11 prints columns 11 (left) and 12 (right); column 12 begins at ``14. The device of claim 13''.}
\textbf{12}. The device of claim \textbf{11} further including a vibrating blade support rail secured to said support arm and supporting said vibrating blade support carriage.

\textbf{13}. The device of claim \textbf{3} further including a blade holder on said vibrating blade support.

\textbf{14}. The device of claim \textbf{13} further including a horizontal ridge on said blade holder and adapted to define a location for the attachment of a cutting blade to said blade holder.

\end{hangpars}

\begin{center}
* * * * *
\end{center}

\end{document}
